\section{Formulasi \textit{Exact Potential Game} (EPG) dan Pencarian \textit{Nash Equilibrium}}
\label{sec:4_2_epg_nash}

Implementasi kerangka kerja \textit{Game Theory} dalam penelitian ini dimulai dengan memetakan permasalahan optimasi Markowitz ke dalam domain \textit{Exact Potential Game} (EPG). Dalam sistem ini, setiap aset dianggap sebagai agen strategis yang berusaha memaksimalkan fungsi utilitas individunya melalui pemilihan status kepemilikan biner $x_i \in \{0, 1\}$. Berdasarkan teori yang dikembangkan oleh Monderer dan Shapley, sebuah permainan dikategorikan sebagai EPG apabila terdapat sebuah fungsi potensial skalar $\Phi(\vec{x})$ sedemikian rupa sehingga selisih utilitas individu $\Delta u_i$ akibat perubahan strategi satu agen secara eksak sama dengan selisih nilai fungsi potensial tersebut. Konstruksi fungsi potensial untuk portofolio Markowitz didefinisikan sebagai berikut:
\begin{equation}
    \Phi(\vec{x}) = \sum_{l=1}^N \mu_l x_l - \frac{\lambda}{2} \sum_{i=1}^N \sum_{j=1}^N \sigma_{ij} x_i x_j
\end{equation}
di mana $\mu_l$ adalah imbal hasil ekspektasi, $\sigma_{ij}$ adalah kovarians antar aset, dan $\lambda$ adalah parameter \textit{risk aversion}.

Eksistensi fungsi potensial $\Phi(\vec{x})$ memberikan jaminan matematis terhadap keberadaan setidaknya satu \textit{Pure Strategy Nash Equilibrium} (PSNE), yang merupakan titik di mana tidak ada agen yang dapat meningkatkan utilitasnya dengan mengubah strateginya secara unilateral. Untuk menemukan titik equilibrium tersebut secara numerik, penelitian ini menggunakan algoritma \textit{Stochastic Best Response} (SBR). Dinamika pencarian equilibrium ini tercatat dalam \textit{file} \texttt{riwayat\_nash\_sbr\_N2.csv}, yang menunjukkan transisi strategi agen dari kondisi awal hingga mencapai konvergensi stabil. Sebagai contoh, pada jendela waktu tertentu, sistem dapat mengalami transisi status dari $|10\rangle$ menuju $|01\rangle$ apabila utilitas marginal yang dihasilkan oleh aset kedua lebih besar daripada aset pertama setelah mempertimbangkan faktor risiko dan korelasi pasar.

Pembuktian ekuivalensi antara selisih utilitas individu dengan selisih fungsi potensial dapat ditunjukkan melalui derivasi perubahan status agen ke-$i$ dari $x_i$ menjadi $x_i^\prime$:
\begin{equation}
    \Delta \Phi = \Phi(x_i, \vec{x}_{-i}) - \Phi(x_i^{\prime}, \vec{x}_{-i}) = u_i(x_i, \vec{x}_{-i}) - u_i(x_i^{\prime}, \vec{x}_{-i})
\end{equation}
Hubungan pada Persamaan (4.9) memastikan bahwa setiap langkah peningkatan utilitas individu yang dilakukan oleh agen secara otomatis akan menurunkan total energi potensial sistem. Oleh karena itu, pencarian PSNE melalui mekanisme SBR bukan hanya sekadar proses heuristik, melainkan sebuah metode yang konsisten secara matematis untuk menemukan konfigurasi portofolio yang paling stabil secara strategis. Hasil akhir dari tahapan ini memberikan solusi klasik definitif yang akan berperan krusial dalam memandu proses optimasi kuantum pada tahapan selanjutnya.
